\documentclass[11pt]{article}

\usepackage[a4paper,margin=1in]{geometry}
\usepackage{amsmath,amssymb,amsthm,mathtools}
\usepackage{mathrsfs}
\usepackage{hyperref}
\usepackage{enumitem}
\usepackage{microtype}

\hypersetup{
	colorlinks=true,
	linkcolor=blue,
	urlcolor=blue,
	citecolor=blue
}

\title{Lecture Notes: Counting \(k\)-Dimensional Subspaces of an \(n\)-Dimensional Vector Space}
\author{}
\date{}

\newtheorem{theorem}{Theorem}[section]
\newtheorem{proposition}[theorem]{Proposition}
\newtheorem{lemma}[theorem]{Lemma}
\newtheorem{corollary}[theorem]{Corollary}
\theoremstyle{definition}
\newtheorem{definition}[theorem]{Definition}
\newtheorem{example}[theorem]{Example}
\newtheorem{exercise}[theorem]{Exercise}
\theoremstyle{remark}
\newtheorem{remark}[theorem]{Remark}

\newcommand{\F}{\mathbf{F}}
\newcommand{\R}{\mathbf{R}}
\newcommand{\C}{\mathbf{C}}
\newcommand{\Q}{\mathbf{Q}}
\newcommand{\Gr}{\mathrm{Gr}}
\newcommand{\GL}{\mathrm{GL}}
\newcommand{\Span}{\operatorname{span}}
\newcommand{\qbinom}[2]{\genfrac{[}{]}{0pt}{}{#1}{#2}_q}

\begin{document}
	
	\maketitle
	
	\tableofcontents
	
	\section{The Main Question}
	
	Let \(V\) be an \(n\)-dimensional vector space. We want to answer the following question:
	
	\begin{quote}
		How many distinct \(k\)-dimensional subspaces does \(V\) have?
	\end{quote}
	
	The answer depends on the field over which \(V\) is defined.
	
	\begin{itemize}
		\item If the field is finite, say \(\F_q\) with \(q\) elements, then the answer is finite and is given by the \emph{Gaussian binomial coefficient}.
		\item If the field is infinite, such as \(\R\), \(\C\), or \(\Q\), then there are infinitely many \(k\)-dimensional subspaces whenever \(0<k<n\).
	\end{itemize}
	
	These notes explain the finite-field formula, prove it carefully, and place it in a broader geometric and combinatorial context.
	
	\section{Ordinary Binomial Coefficients as Motivation}
	
	Recall that the number of \(k\)-element subsets of an \(n\)-element set is
	\[
	\binom{n}{k}.
	\]
	The Gaussian binomial coefficient is the vector-space analogue of this number.
	
	\begin{center}
		\begin{tabular}{c|c}
			Set Theory & Linear Algebra over \(\F_q\) \\ \hline
			\(k\)-element subsets of an \(n\)-element set & \(k\)-dimensional subspaces of \(\F_q^n\) \\
			\(\binom{n}{k}\) & \(\qbinom{n}{k}\)
		\end{tabular}
	\end{center}
	
	So the Gaussian binomial coefficient is often called a \(q\)-analogue of the ordinary binomial coefficient.
	
	\section{Statement of the Formula over a Finite Field}
	
	Assume from now on that
	\[
	V \cong \F_q^n,
	\]
	where \(\F_q\) is a finite field with \(q\) elements.
	
	\begin{definition}
		The \emph{Gaussian binomial coefficient} (or \emph{\(q\)-binomial coefficient}) is defined by
		\[
		\qbinom{n}{k}
		=
		\frac{(q^n-1)(q^n-q)(q^n-q^2)\cdots(q^n-q^{k-1})}
		{(q^k-1)(q^k-q)(q^k-q^2)\cdots(q^k-q^{k-1})}
		\]
		for \(0\leq k\leq n\).
	\end{definition}
	
	An equivalent product form is
	\[
	\qbinom{n}{k}
	=
	\prod_{i=0}^{k-1}\frac{q^n-q^i}{q^k-q^i}.
	\]
	
	Another useful form is
	\[
	\qbinom{n}{k}
	=
	\frac{(q^n-1)(q^{n-1}-1)\cdots(q^{n-k+1}-1)}
	{(q^k-1)(q^{k-1}-1)\cdots(q-1)}.
	\]
	
	\begin{theorem}
		Let \(V\) be an \(n\)-dimensional vector space over \(\F_q\). Then the number of distinct \(k\)-dimensional subspaces of \(V\) is
		\[
		\boxed{
			\qbinom{n}{k}
			=
			\prod_{i=0}^{k-1}\frac{q^n-q^i}{q^k-q^i}
		}.
		\]
	\end{theorem}
	
	\section{Proof by Counting Ordered Bases}
	
	The proof is a clean double-counting argument.
	
	\subsection{Counting Ordered Linearly Independent \(k\)-Tuples in \(\F_q^n\)}
	
	We count the number of ordered tuples
	\[
	(v_1,\dots,v_k)\in (\F_q^n)^k
	\]
	such that \(v_1,\dots,v_k\) are linearly independent.
	
	\begin{itemize}
		\item The first vector \(v_1\) can be any nonzero vector, so there are \(q^n-1\) choices.
		\item The second vector \(v_2\) must not lie in \(\Span(v_1)\), and \(\Span(v_1)\) has \(q\) elements, so there are \(q^n-q\) choices.
		\item The third vector \(v_3\) must not lie in \(\Span(v_1,v_2)\), which has \(q^2\) elements, so there are \(q^n-q^2\) choices.
	\end{itemize}
	
	Continuing in this way, we obtain:
	
	\begin{lemma}
		The number of ordered linearly independent \(k\)-tuples in \(\F_q^n\) is
		\[
		(q^n-1)(q^n-q)(q^n-q^2)\cdots(q^n-q^{k-1}).
		\]
	\end{lemma}
	
	\subsection{Counting Ordered Bases of a Fixed \(k\)-Dimensional Subspace}
	
	Now fix a \(k\)-dimensional subspace \(W\subseteq \F_q^n\). Since \(W\cong \F_q^k\), the number of ordered bases of \(W\) is the same as the number of ordered linearly independent \(k\)-tuples in \(\F_q^k\).
	
	\begin{lemma}
		A fixed \(k\)-dimensional subspace \(W\) has exactly
		\[
		(q^k-1)(q^k-q)(q^k-q^2)\cdots(q^k-q^{k-1})
		\]
		ordered bases.
	\end{lemma}
	
	\subsection{Dividing Out the Overcount}
	
	Each ordered linearly independent \(k\)-tuple spans a unique \(k\)-dimensional subspace, and every \(k\)-dimensional subspace contributes exactly the number of ordered bases computed above.
	
	Therefore:
	
	\begin{proof}[Proof of the Theorem]
		The number of ordered linearly independent \(k\)-tuples in \(\F_q^n\) is
		\[
		(q^n-1)(q^n-q)\cdots(q^n-q^{k-1}).
		\]
		Each \(k\)-dimensional subspace has exactly
		\[
		(q^k-1)(q^k-q)\cdots(q^k-q^{k-1})
		\]
		ordered bases. Dividing, we obtain the number of distinct \(k\)-dimensional subspaces:
		\[
		\frac{(q^n-1)(q^n-q)\cdots(q^n-q^{k-1})}
		{(q^k-1)(q^k-q)\cdots(q^k-q^{k-1})}
		=
		\qbinom{n}{k}.
		\]
	\end{proof}
	
	\section{Why the Denominator Appears}
	
	A frequent source of confusion is the denominator
	\[
	(q^k-1)(q^k-q)\cdots(q^k-q^{k-1}).
	\]
	
	The reason is that many different ordered bases span the same subspace. So the numerator counts the same subspace many times, once for each ordered basis of that subspace.
	
	Thus:
	\begin{itemize}
		\item numerator = all ordered bases of all \(k\)-dimensional subspaces,
		\item denominator = ordered bases of one fixed \(k\)-dimensional subspace,
		\item quotient = number of distinct \(k\)-dimensional subspaces.
	\end{itemize}
	
	\section{Examples}
	
	\begin{example}[Lines in \(\F_q^3\)]
		A \(1\)-dimensional subspace is a line through the origin. The number of such lines is
		\[
		\qbinom{3}{1}
		=
		\frac{q^3-1}{q-1}
		=
		q^2+q+1.
		\]
	\end{example}
	
	\begin{example}[Planes in \(\F_q^3\)]
		A \(2\)-dimensional subspace of \(\F_q^3\) is a plane through the origin. The number of such planes is
		\[
		\qbinom{3}{2}
		=
		\qbinom{3}{1}
		=
		q^2+q+1.
		\]
	\end{example}
	
	\begin{example}[Two-dimensional subspaces of \(\F_2^4\)]
		Take \(q=2\), \(n=4\), \(k=2\). Then
		\[
		\qbinom{4}{2}
		=
		\frac{(2^4-1)(2^4-2)}{(2^2-1)(2^2-2)}
		=
		\frac{15\cdot 14}{3\cdot 2}
		=
		35.
		\]
		So \(\F_2^4\) has exactly \(35\) distinct \(2\)-dimensional subspaces.
	\end{example}
	
	\begin{example}[All lines in \(\F_2^2\)]
		The vector space \(\F_2^2\) has vectors
		\[
		(0,0),\quad (1,0),\quad (0,1),\quad (1,1).
		\]
		Its nonzero vectors are
		\[
		(1,0),\quad (0,1),\quad (1,1).
		\]
		Over \(\F_2\), each \(1\)-dimensional subspace has exactly one nonzero vector, so the distinct lines are
		\[
		\Span(1,0),\qquad \Span(0,1),\qquad \Span(1,1).
		\]
		Thus there are \(3\) of them, in agreement with
		\[
		\qbinom{2}{1}=\frac{2^2-1}{2-1}=3.
		\]
	\end{example}
	
	\section{Special Cases}
	
	\begin{proposition}
		For all \(n\) and \(q\), we have:
		\begin{align*}
			\qbinom{n}{0} &= 1,\\
			\qbinom{n}{1} &= \frac{q^n-1}{q-1}=1+q+q^2+\cdots+q^{n-1},\\
			\qbinom{n}{n-1} &= \qbinom{n}{1},\\
			\qbinom{n}{n} &= 1.
		\end{align*}
	\end{proposition}
	
	\begin{proof}
		The cases \(k=0\) and \(k=n\) correspond to the unique zero subspace and the whole space. The case \(k=1\) is immediate from the formula. The case \(k=n-1\) follows from symmetry, proved below.
	\end{proof}
	
	\section{Symmetry}
	
	Just as ordinary binomial coefficients satisfy
	\[
	\binom{n}{k}=\binom{n}{n-k},
	\]
	Gaussian binomial coefficients satisfy the analogous identity.
	
	\begin{proposition}
		For \(0\leq k\leq n\),
		\[
		\qbinom{n}{k}=\qbinom{n}{n-k}.
		\]
	\end{proposition}
	
	\begin{proof}
		Using the form
		\[
		\qbinom{n}{k}
		=
		\frac{(q^n-1)(q^{n-1}-1)\cdots(q^{n-k+1}-1)}
		{(q^k-1)(q^{k-1}-1)\cdots(q-1)},
		\]
		one checks directly that replacing \(k\) by \(n-k\) produces the same value after cancellation.
	\end{proof}
	
	\begin{remark}
		This symmetry reflects a geometric duality: a \(k\)-dimensional subspace is naturally related to an \((n-k)\)-dimensional object in the dual vector space.
	\end{remark}
	
	\section{A Recursive Formula}
	
	Ordinary binomial coefficients satisfy Pascal's identity
	\[
	\binom{n}{k}=\binom{n-1}{k}+\binom{n-1}{k-1}.
	\]
	Gaussian binomial coefficients satisfy a \(q\)-analogue.
	
	\begin{proposition}
		For \(1\leq k\leq n-1\),
		\[
		\qbinom{n}{k}
		=
		\qbinom{n-1}{k}
		+
		q^{\,n-k}\qbinom{n-1}{k-1}.
		\]
	\end{proposition}
	
	\begin{proof}
		Fix a hyperplane \(H\subseteq \F_q^n\) of dimension \(n-1\). Let \(W\subseteq \F_q^n\) be a \(k\)-dimensional subspace.
		
		There are two cases.
		
		\begin{enumerate}[label=\arabic*.]
			\item \(W\subseteq H\).  
			The number of such subspaces is \(\qbinom{n-1}{k}\).
			
			\item \(W\not\subseteq H\).  
			Then \(W\cap H\) has dimension \(k-1\). For each \((k-1)\)-dimensional subspace \(U\subseteq H\), the number of \(k\)-dimensional subspaces \(W\) such that \(W\cap H=U\) is \(q^{n-k}\). Hence the number of subspaces in this case is
			\[
			q^{n-k}\qbinom{n-1}{k-1}.
			\]
		\end{enumerate}
		
		Adding the two cases gives
		\[
		\qbinom{n}{k}
		=
		\qbinom{n-1}{k}
		+
		q^{n-k}\qbinom{n-1}{k-1}.
		\]
	\end{proof}
	
	\section{Polynomial Nature}
	
	Although the formula for \(\qbinom{n}{k}\) is written as a fraction, it actually simplifies to a polynomial in \(q\) with integer coefficients.
	
	\begin{example}
		\[
		\qbinom{3}{1}=q^2+q+1.
		\]
		Also,
		\[
		\qbinom{4}{2}=q^4+q^3+2q^2+q+1.
		\]
	\end{example}
	
	This fact is highly important in algebraic combinatorics.
	
	\section{The \(q\to 1\) Philosophy}
	
	The Gaussian binomial coefficient is called a \(q\)-analogue of \(\binom{n}{k}\) because, after simplification as a polynomial in \(q\), substituting \(q=1\) gives the ordinary binomial coefficient.
	
	\begin{example}
		Since
		\[
		\qbinom{3}{1}=q^2+q+1,
		\]
		setting \(q=1\) gives
		\[
		1+1+1=3=\binom{3}{1}.
		\]
	\end{example}
	
	So one may think of
	\[
	\qbinom{n}{k}
	\]
	as the vector-space version of
	\[
	\binom{n}{k}.
	\]
	
	\section{What Happens over Infinite Fields?}
	
	Suppose now that the base field is infinite.
	
	\begin{proposition}
		If \(V\) is an \(n\)-dimensional vector space over an infinite field and \(0<k<n\), then \(V\) has infinitely many distinct \(k\)-dimensional subspaces.
	\end{proposition}
	
	\begin{proof}
		It is enough to consider the case \(k=1\). In \(\R^2\), each line through the origin is a \(1\)-dimensional subspace, and such lines are parametrized by slope:
		\[
		y=mx \qquad (m\in \R),
		\]
		together with the vertical line \(x=0\). Thus there are infinitely many \(1\)-dimensional subspaces. The general case follows similarly.
	\end{proof}
	
	\begin{remark}
		The only finite cases are trivial:
		\begin{itemize}
			\item exactly one \(0\)-dimensional subspace, namely \(\{0\}\),
			\item exactly one \(n\)-dimensional subspace, namely \(V\).
		\end{itemize}
	\end{remark}
	
	\section{The Grassmannian Viewpoint}
	
	\begin{definition}
		The set of all \(k\)-dimensional subspaces of an \(n\)-dimensional vector space is called the \emph{Grassmannian}, denoted
		\[
		\Gr(k,n).
		\]
	\end{definition}
	
	Over a finite field, the Grassmannian is a finite set, and
	\[
	|\Gr(k,n)(\F_q)|=\qbinom{n}{k}.
	\]
	
	Over \(\R\) or \(\C\), the Grassmannian is not finite. Instead, it is a geometric object: a manifold and also an algebraic variety.
	
	So the Gaussian binomial coefficient counts the \(\F_q\)-rational points of the Grassmannian.
	
	\section{A Matrix-Theoretic Interpretation}
	
	There is another very useful proof using matrices.
	
	A \(k\)-dimensional subspace of \(\F_q^n\) can be represented as the row space of a full-rank \(k\times n\) matrix.
	
	\begin{itemize}
		\item The number of full-rank \(k\times n\) matrices is
		\[
		(q^n-1)(q^n-q)\cdots(q^n-q^{k-1}).
		\]
		\item Two such matrices have the same row space if and only if one is obtained from the other by multiplication on the left by an invertible \(k\times k\) matrix.
		\item The group \(\GL_k(\F_q)\) has size
		\[
		|\GL_k(\F_q)|=(q^k-1)(q^k-q)\cdots(q^k-q^{k-1}).
		\]
	\end{itemize}
	
	Therefore,
	\[
	\#\{\text{\(k\)-dimensional subspaces of }\F_q^n\}
	=
	\frac{\#\{\text{full-rank }k\times n\text{ matrices}\}}
	{|\GL_k(\F_q)|}
	=
	\qbinom{n}{k}.
	\]
	
	\section{Related Formula: The Size of \(\GL_n(\F_q)\)}
	
	The same counting method gives the size of the general linear group.
	
	\begin{proposition}
		\[
		|\GL_n(\F_q)|=(q^n-1)(q^n-q)(q^n-q^2)\cdots(q^n-q^{n-1}).
		\]
	\end{proposition}
	
	\begin{proof}
		An invertible \(n\times n\) matrix is exactly an ordered basis of \(\F_q^n\). The first column can be any nonzero vector, the second must be outside the span of the first, and so on.
	\end{proof}
	
	This is the full-basis version of the same principle used in the Gaussian binomial coefficient.
	
	\section{Conceptual Summary}
	
	The formula
	\[
	\qbinom{n}{k}
	=
	\frac{(q^n-1)(q^n-q)\cdots(q^n-q^{k-1})}
	{(q^k-1)(q^k-q)\cdots(q^k-q^{k-1})}
	\]
	comes from two ideas:
	
	\begin{enumerate}[label=\arabic*.]
		\item Count ordered linearly independent \(k\)-tuples in \(\F_q^n\).
		\item Divide by the number of ordered bases of a fixed \(k\)-dimensional subspace.
	\end{enumerate}
	
	This is the exact finite-field analogue of counting \(k\)-element subsets of an \(n\)-element set.
	
	\section{Final Theorem}
	
	\begin{theorem}[Final Form]
		Let \(V\) be an \(n\)-dimensional vector space over \(\F_q\). Then the number of distinct \(k\)-dimensional subspaces of \(V\) is
		\[
		\boxed{
			\qbinom{n}{k}
			=
			\prod_{i=0}^{k-1}\frac{q^n-q^i}{q^k-q^i}
		}.
		\]
		If the base field is infinite, then for \(0<k<n\) the number of \(k\)-dimensional subspaces is infinite.
	\end{theorem}
	
	\section{Exercises}
	
	\begin{exercise}
		Compute \(\qbinom{5}{1}\) when \(q=2\).
	\end{exercise}
	
	\begin{exercise}
		Compute \(\qbinom{4}{2}\) when \(q=3\).
	\end{exercise}
	
	\begin{exercise}
		Show directly that
		\[
		\qbinom{n}{1}=1+q+q^2+\cdots+q^{n-1}.
		\]
	\end{exercise}
	
	\begin{exercise}
		Prove that
		\[
		\qbinom{n}{k}=\qbinom{n}{n-k}.
		\]
	\end{exercise}
	
	\begin{exercise}
		Verify the recursion
		\[
		\qbinom{n}{k}
		=
		\qbinom{n-1}{k}
		+
		q^{n-k}\qbinom{n-1}{k-1}.
		\]
	\end{exercise}
	
	\begin{exercise}
		Count the \(2\)-dimensional subspaces of \(\F_2^3\).
	\end{exercise}
	
	\begin{exercise}
		Explain why there are infinitely many \(1\)-dimensional subspaces of \(\R^2\).
	\end{exercise}
	
	\section{Selected Answers}
	
	\begin{itemize}
		\item For \(q=2\),
		\[
		\qbinom{5}{1}=\frac{2^5-1}{2-1}=31.
		\]
		
		\item For \(q=3\),
		\[
		\qbinom{4}{2}
		=
		\frac{(3^4-1)(3^4-3)}{(3^2-1)(3^2-3)}
		=
		\frac{80\cdot 78}{8\cdot 6}
		=
		130.
		\]
		
		\item
		\[
		\binom{3}{2}_{q=2}
		=
		\binom{3}{1}_{q=2}
		=
		\frac{2^3-1}{2-1}
		=
		7.
		\]
	\end{itemize}
	
	\section*{One-Sentence Takeaway}
	
	Over a finite field \(\F_q\), the number of \(k\)-dimensional subspaces of \(\F_q^n\) is the Gaussian binomial coefficient \(\qbinom{n}{k}\); over an infinite field, the number is usually infinite.
	
\end{document}